\section{Fixed-packet asymptotics of the complete original Gamma jet determinant}
\label{sec:fixed-gamma-full-jet}

Fix the original positive integer $k$, put $c=k/2$, and retain the
actual monic cyclic annihilator $\chi=\chi_{h,k}$, with all its
distinct roots $\lambda_1,\ldots,\lambda_d$ and their complete
multiplicities $m_1,\ldots,m_d$. Write $q=\sum_i m_i$.
The finite-packet hypothesis only specifies these original objects;
no new arithmetic zero or replacement annihilator enters this section.
The same computation applies to every specified finite jet observation
on the given source line. This last fact will be used to transport the
original observation through a literal comparison-measure dilation.

The source line and comparison density are exactly
\[
 S=c+iu,\qquad
 \sigma(u)=\frac{|\Gamma(1/4+iu/2)|^2}{2\pi}.
 \tag{FPK.1}
\]
Use the original monic polynomials and squared norms
\[
 p_n(S)=i^n b_n((S-c)/i),\qquad
 \gamma_n=\sqrt{2\pi}\,n!(1/2)_n
          =\sqrt2\,n!\Gamma(n+1/2).
 \tag{FPK.2}
\]
The generating-integral calculation in the retained Gamma chapter proves
their orthogonality for (FPK.1), including these constants and phases.
We calculate the full confluent kernel, not just the diagonal scalar
evaluation kernel.

For $f^{[j]}=f^{(j)}/j!$, let $\mathsf E_\chi$ be the original
Taylor-coordinate map from the increasing-power remainder frame:
\[
 (\mathsf E_\chi)_{(i,j),r}
   =\begin{cases}\binom rj\lambda_i^{r-j},&r\geq j,\\0,&r<j,
     \end{cases}
 \quad 0\leq j<m_i,\quad 0\leq r<q.
 \tag{FPK.3}
\]
Define
\[
 (\mathsf K_N)_{(i,j),(a,l)}
   =\sum_{n=0}^N
      \frac{p_n^{[j]}(\lambda_i)
                   \overline{p_n^{[l]}(\lambda_a)}}{\gamma_n}
       \qquad(N\geq q-1).
 \tag{FPK.4}
\]
The minimum-norm quotient Gram and its original determinant are
\[
 G_N^\sigma=\mathsf E_\chi^*\mathsf K_N^{-1}\mathsf E_\chi,
 \qquad V_N^\sigma=\det G_N^\sigma
         =\frac{|\det\mathsf E_\chi|^2}{\det\mathsf K_N}.
 \tag{FPK.5}
\]
Indeed the normalized polynomial basis makes the source Gram identity,
so the Gram of the onto jet observation is (FPK.4). Minimizing a
Euclidean norm subject to that observation gives its inverse Gram.
The first $q$ monic polynomials span all remainders, proving that this
Gram is positive definite. Transport by the actual map (FPK.3) gives
(FPK.5). Thus $V_N$ means the determinant of the Gram, as in (AT6),
and does not mean its square root.

The exact determinant of the displayed coordinate map is
\[
 \det\mathsf E_\chi
       =\prod_{i<a}(\lambda_a-\lambda_i)^{m_i m_a}.
 \tag{FPK.6}
\]
Here root blocks and derivatives are ordered as in (FPK.3).
To prove the constant, replace the $m_i$ evaluations in block $i$
by evaluations at $\lambda_i+\epsilon x_{i,j}$ with fixed distinct
$x_{i,j}$. Expand each row into its Taylor series and divide its
block determinant by
$\epsilon^{m_i(m_i-1)/2}\prod_{j<l}(x_{i,l}-x_{i,j})$.
Multilinearity gives the determinant of the normalized derivative
rows in the limit. The ordinary Vandermonde formula on all the
evaluation points gives precisely the right side of (FPK.6);
the within-block factors cancel the displayed divisors. This also
proves that no factorial or orientation factor is missing.

\subsection{The two singularities with the original coordinate and constants}

Put $w=S-c$. Substituting $iz$ in the retained generating function
$(1+z^2)^{-1/4}\exp(u\arctan z)$ gives the exact identity
\[
 \sum_{n=0}^\infty \frac{p_n(c+w)}{n!}z^n
    =(1-z)^{-a(w)}(1+z)^{-b(w)},
 \quad a(w)=\tfrac14+\tfrac w2,
 \quad b(w)=\tfrac14-\tfrac w2.
 \tag{FPK.7}
\]
All powers use the branches equal to one at $z=0$. Define the entire
coefficient functions
\[
 C_+(w)=\frac{2^{-1/2+w/2}}{\Gamma(1/4+w/2)},\qquad
 C_-(w)=\frac{2^{-1/2-w/2}}{\Gamma(1/4-w/2)}.
 \tag{FPK.8}
\]
For every compact set of $w$ and each fixed derivative order $j$,
the normalized polynomials $e_n(w)=p_n(c+w)/\sqrt{\gamma_n}$ satisfy
\[
 e_n(w)=C_+(w)n^{w/2-1/2}
           +(-1)^n C_-(w)n^{-w/2-1/2}+R_n(w),
 \tag{FPK.9}
\]
\[
 |\partial_w^jR_n(w)|
 \leq C_j(1+\log n)^j
       \left(n^{\Re w/2-3/2}+n^{-\Re w/2-3/2}
                         +n^{-J}\right),\qquad n\geq2.
 \tag{FPK.10}
\]
The positive integer $J$ can be chosen as large as needed, with the
constant depending on the compact set, $j$, and $J$.

Here is a coefficient proof of these estimates. It is included so no
unproved asymptotic theorem about the orthogonal polynomials is used.
Near $z=1$, expand the other factor in (FPK.7):
\[
 (1+z)^{-b}=2^{-b}
   \sum_{r=0}^{L-1}\frac{(b)_r}{2^r r!}(1-z)^r
            +(1-z)^L A_L(z,w),
 \tag{FPK.11}
\]
where $A_L$ is analytic near that point, uniformly with each specified
parameter derivative on a compact parameter set. Near $z=-1$ use the
same expansion with $a,b$ exchanged and $z$ replaced by $-z$.
Subtract from (FPK.7) both finite lists of singular terms. Choose $L$
larger than the maximum real part of $a,b$ plus any specified integer
$J_0+2$. The remainder and its first $J_0$ angular derivatives on the
unit circle then extend continuously across both $1$ and $-1$.
Indeed its nonanalytic local remainder is a power of real part greater
than $J_0+1$, times an analytic bounded factor; parameter derivatives
only add fixed powers of its logarithm. The singular terms subtracted
at the other point are analytic there. Away from these two points
everything is analytic across the circle.

The resulting boundary function is $C^{J_0}$ and periodic. Integrating
its Fourier coefficient by parts $J_0$ times bounds that coefficient
by $C n^{-J_0}$, uniformly for all the specified parameter derivatives.
Its nonnegative Fourier coefficients are its Taylor coefficients:
apply Cauchy's formula on circles of radius less than one and use
uniform boundary convergence. Each subtracted term has its exact
coefficient
\[
 [z^n](1-z)^{-a+r}
   =\frac{\Gamma(n+a-r)}{\Gamma(a-r)\Gamma(n+1)},
 \tag{FPK.12}
\]
with analytic values at nonpositive integral $a-r$; the corresponding
coefficient at $-1$ has the additional factor $(-1)^n$.

For completeness, uniformly on a compact complex set of $\alpha$,
\[
 \frac{\Gamma(n+\alpha)}{\Gamma(n+1)}
       =n^{\alpha-1}(1+O(n^{-1})),
 \tag{FPK.13}
\]
with differentiated estimates after retaining the derivatives of
$n^{\alpha-1}$. One direct proof first arranges
$\Re\alpha<1$ by the gamma recurrence. The beta integral is
\[
 \frac{\Gamma(n+\alpha)}{\Gamma(n+1)}
   =\frac{n^{\alpha-1}}{\Gamma(1-\alpha)}
     \int_0^n(1-y/n)^{n+\alpha-1}y^{-\alpha}\,dy.
\]
On $0\leq y\leq n/2$, the Taylor estimate for $\log(1-y/n)$,
with its integral remainder, bounds the difference from $e^{-y}$
in the integral by $C n^{-1}e^{-y/2}(1+y)^2y^{-\Re\alpha}$.
One may first use $0\leq y\leq n/4$ for this estimate and absorb
$n/4\leq y\leq n/2$ into its exponentially small bound. On
$n/2\leq y\leq n$, integrate the exponentially small factor
$(1-y/n)^{n+\Re\alpha-1}$ directly. The same estimates hold after
each fixed parameter derivative, with additional integrable powers
of $|\log y|$ and $|\log(1-y/n)|$. The limiting integral is
$\Gamma(1-\alpha)$; its reciprocal is bounded on the chosen compact
set. This proves (FPK.13) and its parameter derivatives. Recurrence
back to the original $\alpha$ multiplies by a fixed finite product
of $n+\alpha-r$, yielding the asserted formula on every compact set.

In (FPK.12), the $r=0$ terms therefore have coefficients
\[
 2^{-b}n^{a-1}/\Gamma(a),\qquad
 (-1)^n2^{-a}n^{b-1}/\Gamma(b),
\]
each with a relative error of order
$1/n$ interpreted analytically when a reciprocal gamma factor
vanishes. The terms $r\geq1$ are of the next indicated order or
smaller. The exact normalization (FPK.2) multiplies these coefficients
by
\[
 \frac{n!}{\sqrt{\gamma_n}}
    =2^{-1/4}\left(\frac{\Gamma(n+1)}{\Gamma(n+1/2)}\right)^{1/2}
    =2^{-1/4}n^{1/4}(1+O(n^{-1})).
\]
This proves (FPK.8)--(FPK.10), including the original factors of two
and every derivative required for the full jet calculation.

\subsection{The complete off-critical blocks}

For $w_i=\lambda_i-c$, put $\delta_i=\Re w_i$. When
$\delta_i\ne0$, set $\epsilon_i=\operatorname{sgn}\delta_i$.
The dominant coefficient $C_{\epsilon_i}(w)$ is nonzero on a
neighborhood of $w_i$, because the real part of its gamma argument
at $w_i$ is $1/4+|\delta_i|/2>0$. On that full local jet block
define $\mathsf T_{N,i}$ to be multiplication by the analytic germ
\[
 \left(C_{\epsilon_i}(w)N^{\epsilon_i w/2}\right)^{-1},
 \qquad \det\mathsf T_{N,i}
       =\left(C_{\epsilon_i}(w_i)N^{\epsilon_iw_i/2}\right)^{-m_i}.
 \tag{FPK.14}
\]
This retains all the derivatives of $C_{\epsilon_i}$ and the
triangular logarithmic jet terms. It is a specified invertible
coordinate map used to compute the determinant, and will be undone.

Apply it to the complete Taylor column of $e_n$ at $w_i$. The
dominant component in position $j$ becomes
\[
 v_{N,n;i,j}
  =\vartheta_{i,n}\,n^{-1/2}(n/N)^{\epsilon_iw_i/2}
             \frac{(\epsilon_i\log(n/N)/2)^j}{j!},
 \quad
 \vartheta_{i,n}=\begin{cases}1,&\epsilon_i=1,\\(-1)^n,&\epsilon_i=-1.
 \end{cases}
 \tag{FPK.15}
\]
The difference between the normalized exact row and (FPK.15) has
squared-sum norm tending to zero. To verify that
claim without neglecting small indices, the nondominant term and
each of its fixed derivatives have squared-sum bound
\[
 C N^{-|\delta_i|}(1+\log N)^{2j}
       \sum_{n=1}^N n^{-1-|\delta_i|}(1+\log n)^{2j}=o(1).
\]
For the remainder (FPK.10), multiplication and differentiation in
(FPK.14) give a bound by a constant times a fixed power of $\log N$
times
\[
 N^{-|\delta_i|}\sum_{n=2}^N n^{|\delta_i|-3}
       +N^{-|\delta_i|}\sum_{n=2}^N n^{-1-|\delta_i|}.
\]
It tends to zero: the first term is $O(N^{-2})$ if
$|\delta_i|>2$, $O(N^{-2}\log N)$ if it equals two, and
$O(N^{-|\delta_i|})$ otherwise. Choose $J$ in (FPK.10) sufficiently
large for the remaining rapidly decaying term. The finitely many
initial columns obey the same vanishing conclusion directly, because
their normalization is $O(N^{-|\delta_i|/2}(\log N)^j)$.

For two blocks in the same half-plane, the Gram of (FPK.15) converges
by an integrable Riemann sum to
\[
 (\mathsf H_\epsilon)_{(i,j),(a,l)}
  =\int_0^1 x^{\epsilon(w_i+\overline w_a)/2-1}
       \frac{(\epsilon\log x/2)^{j+l}}{j!l!}\,dx
  =\frac{(-\epsilon/2)^{j+l}(j+l)!}
      {j!l!\,[\epsilon(w_i+\overline w_a)/2]^{j+l+1}}.
 \tag{FPK.16}
\]
Here $\epsilon_i=\epsilon_a=\epsilon$, and the denominator has
positive real part. Convergence at zero follows, for every fixed
power of $\log x$, from that positive real part. For a direct
Riemann-sum justification, split at any $b>0$; on $[b,1]$ the
integrand is smooth, and on $(0,b)$ its absolute sums and integrals
are bounded by $C b^\alpha(1+|\log b|^{j+l})$, with some
$\alpha>0$, which tends to zero.

The matrix $\mathsf H_\epsilon$ is positive definite. It is the
Gram of the complete list
$x^{\epsilon w_i/2-1/2}(\epsilon\log x/2)^j/j!$ in $L^2(0,1)$.
An equality between these functions, after $x=e^{-y}$ and division
by the common nonzero factor $e^{y/2}$, gives an
exponential-polynomial equality on $y>0$. Apply the product of
$(\partial_y+\epsilon w_a/2)^{m_a}$ over $a\ne i$.
It removes the other blocks and acts on the remaining polynomial
by a product of nonzero constants plus nilpotent differentiation.
That product is invertible on the specified finite polynomial space.
Every remaining polynomial is therefore zero, proving independence.

For two opposite half-planes the Gram has instead an additional
$(-1)^n$. Its absolute integrable Riemann envelope is the same
with real exponent $(|\delta_i|+|\delta_a|)/2>0$ at zero.
On $[b,1]$, pairing successive indices proves that its alternating
Riemann sum tends to zero by uniform continuity. The absolute
bound on $(0,b)$ just proved then gives a zero limit for the whole
sum. Thus the two half-plane blocks decouple in this exact limit;
no entry of the original finite kernel has been set to zero.

\subsection{Critical full jets and all mixed blocks}

For $\delta_i=0$, put $L_N=\log N$ and use the diagonal jet map
\[
 (\mathsf T_{N,i})_{j,j}=L_N^{-j-1/2},
 \qquad \det\mathsf T_{N,i}=L_N^{-m_i^2/2}.
 \tag{FPK.17}
\]
In a differentiated dominant term (FPK.9), the part in which all
$j$ derivatives hit the power of $n$ is
\[
 \frac{(\log n)^j}{2^j j!\,\sqrt n}
       \left(C_+(w_i)n^{w_i/2}
                    +(-1)^{n+j}C_-(w_i)n^{-w_i/2}\right).
 \tag{FPK.18}
\]
Every other differentiated term has a smaller power of $\log n$.
After multiplication by (FPK.17), its squared-sum norm tends to
zero, since $\sum_{n\leq N}n^{-1}(\log n)^r
 =O(1+L_N^{r+1})$. The remainder in (FPK.10) and every fixed
initial column vanish in the same normalized squared sum.

At the same critical point, the nonalternating products in (FPK.18)
give the limit
\[
 (\mathsf H_i^0)_{j,l}
 =\frac{|C_+(w_i)|^2+(-1)^{j+l}|C_-(w_i)|^2}
            {2^{j+l}j!l!(j+l+1)}
 =\frac{|C_+(w_i)|^2}{2^{j+l}j!l!}
                    \int_{-1}^1t^{j+l}\,dt.
 \tag{FPK.19}
\]
The last equality uses the exact relation
$C_-(w_i)=\overline{C_+(w_i)}$ for purely imaginary $w_i$.
The coefficient is nonzero, and (FPK.19) is a positive multiple
of the Gram of independent real monomials of degrees below $m_i$.
It is therefore positive definite, including every nilpotent direction.

For different critical points, a nonalternating product has
$\sum_{n\leq N}n^{-1+i\theta}(\log n)^r$ with $\theta\ne0$.
Its difference from the corresponding integral is bounded uniformly
in $N$: the derivative of $x^{-1+i\theta}(\log x)^r$ is absolutely
integrable on $[1,\infty)$. Substitution $x=e^t$ and integration
by parts bound that integral by $O(1+L_N^r)$. Division by
$L_N^{r+1}$ gives zero. Alternating products are bounded uniformly
before this division, by pairing successive terms and the same
integrable derivative estimate. This includes the alternating
products within a single critical block. Hence all different critical
blocks have zero cross limit.

Every mixed off-critical/critical block also has zero limit. Here
is a bound that does not require any oscillatory cancellation.
The normalized off-critical row (FPK.15) has a limiting squared
mass distribution on $x=n/N\in(0,1)$ with integrable density
$x^{|\delta_i|-1}$ times a fixed power of $\log x$.
For any fixed $R>0$, its mass in $n\leq Ne^{-R}$ has limiting
upper bound equal to the integral of that density on $(0,e^{-R})$,
which tends to zero as $R$ grows. On the complementary indices,
the squared mass of a normalized critical row is $O(R/L_N)$:
use $(\log n)^{2j}\leq L_N^{2j}$ and
$\sum_{Ne^{-R}<n\leq N}1/n\leq R+o(1)$.
Both rows have bounded total squared norm by their already-proved
diagonal limits. Cauchy--Schwarz on the two index ranges, first
letting $N$ tend to infinity and then $R$ tend to infinity, proves
that their inner product tends to zero. The vanishing remainder
rows do not change this conclusion.

\subsection{The entire determinant and the original quotient volume}

Let $\mathsf T_N$ be the block diagonal combination of
(FPK.14) and (FPK.17), with any fixed root-block ordering. The
preceding proofs establish the full finite-matrix limit
\[
 \mathsf T_N\mathsf K_N\mathsf T_N^*
    \longrightarrow
 \mathsf H:=\mathsf H_+\oplus\mathsf H_-
          \oplus\bigoplus_{\delta_i=0}\mathsf H_i^0>0.
 \tag{FPK.20}
\]
Absent blocks are omitted and their determinants are one. Put
\[
 \mathcal A_\chi=\sum_i m_i|\Re\lambda_i-c|,
 \qquad \mathcal B_\chi=\sum_{\Re\lambda_i=c}m_i^2,
 \tag{FPK.21}
\]
\[
 \mathcal C_\chi
  =(\det\mathsf H)
     \prod_{\delta_i\ne0}|C_{\epsilon_i}(\lambda_i-c)|^{2m_i}>0.
 \tag{FPK.22}
\]
All entries defining this constant have been given in (FPK.8),
(FPK.16), and (FPK.19), with their actual coordinates and factors.

\begin{theorem}[Complete fixed-Gamma jet determinant]
\label{thm:fpk-full-jet-asymptotic}
For the original fixed source line, the entire original finite jet
observation, and $N\to\infty$ through integers,
\[
 \det\mathsf K_N
   =\mathcal C_\chi N^{\mathcal A_\chi}
       (\log N)^{\mathcal B_\chi}(1+o(1)),
 \tag{FPK.23}
\]
\[
 V_N^\sigma
   =\frac{|\det\mathsf E_\chi|^2}{\mathcal C_\chi}
       N^{-\mathcal A_\chi}(\log N)^{-\mathcal B_\chi}(1+o(1)),
 \tag{FPK.24}
\]
and therefore
\[
 \log V_N^\sigma
 =-\mathcal A_\chi\log N-\mathcal B_\chi\log\log N
       +\log\frac{|\det\mathsf E_\chi|^2}{\mathcal C_\chi}+o(1).
 \tag{FPK.25}
\]
\end{theorem}
\begin{proof}
Determinant is a polynomial in the matrix entries, so (FPK.20)
gives $\det(\mathsf T_N\mathsf K_N\mathsf T_N^*)
\to\det\mathsf H>0$. The literal determinant factors of the
coordinate maps are
\[
 |\det\mathsf T_N|^2
   =N^{-\mathcal A_\chi}(\log N)^{-\mathcal B_\chi}
        \prod_{\delta_i\ne0}|C_{\epsilon_i}(w_i)|^{-2m_i}.
\]
Undoing this congruence proves (FPK.23). The exact original
coordinate formula (FPK.5) proves (FPK.24). Its positive constant
and positive determinants allow taking real logarithms, which gives
(FPK.25).
\end{proof}

The original arithmetic observation remains
$\sigma_h^{\otimes k}\eta_{h,k}:E\hookrightarrow Q^{\otimes k}$,
with $\eta_{h,k}$ containing its complete Taylor unit
$\upsilon_h^{\otimes k}$. No factor in this injection was changed
in forming the comparison norm on the original remainder coordinate
$E$. If that coordinate is transported by a specified fixed invertible
unit matrix $U$, the quotient Gram becomes $U^*G_N^\sigma U$ and
its determinant is multiplied by the exact fixed factor $|\det U|^2$.
Thus a unit-containing alternative frame retains the same two exponents
and its own explicitly transported constant; the unit is never set to one.

This theorem is a fixed-$k$, fixed-$\chi$ statement. It provides no
uniform error as the arithmetic packet or $k$ changes, and does not
assign a sign to the original signed four-window correction. Its
assertions are exact asymptotic statements for the stated full kernel
and quotient determinant, retaining every original jet.

\subsection{The original tensor-Gamma reference with its full mass}

The original reference in (AT4), on this same line $S=k/2+iu$, is
\[
 m_0(u)=\frac{(2\pi)^{k/2}}{c_{k/4}}
           \frac{|\Gamma(k/4+iu/2)|^2}{2\pi},
 \qquad c_a=2^{1-2a}\Gamma(2a).
 \tag{FPK.26}
\]
Its mass is exactly $(2\pi)^{k/2}$; this scalar is retained.
The original generating-integral calculation gives its monic
polynomials $p_{k,n}(S)=i^n b_n^{(k/4)}((S-c)/i)$ and squared norms
\[
 \gamma_{k,n}=(2\pi)^{k/2}
                  \frac{n!\Gamma(n+k/2)}{\Gamma(k/2)}.
 \tag{FPK.27}
\]
Their generating function in (FPK.7) has $1/4$ replaced by $k/4$.
The coefficient proof (FPK.11)--(FPK.13), with this fixed positive
parameter, gives the normalized expansion (FPK.9)--(FPK.10) with
the exact coefficients
\[
 C_{\pm,k}(w)=
 (2\pi)^{-k/4}\sqrt{\Gamma(k/2)}\,
          \frac{2^{-k/4\pm w/2}}{\Gamma(k/4\pm w/2)}.
 \tag{FPK.28}
\]
To verify the unchanged power of $n$, a singular coefficient is
$n^{k/4\pm w/2-1}$, while division of $n!$ by the square root
of (FPK.27) contributes $n^{1/2-k/4}$ and the displayed prefactor.
Their product is $n^{\pm w/2-1/2}$. The next-order coefficient
and every derivative have exactly the bounds already used above.
The dominant coefficients remain nonzero, since their gamma
arguments have positive real part. At critical points the two
coefficients are conjugate. Thus every full-block estimate,
positive Gram limit, and determinant argument (FPK.14)--(FPK.25)
applies with (FPK.28), keeping the original root coordinates and
the same $\mathsf E_\chi$. In particular
\[
 \log V_N^{m_0}
   =-\mathcal A_\chi\log N-\mathcal B_\chi\log\log N
       +\log\frac{|\det\mathsf E_\chi|^2}{\mathcal C_{\chi,k}}+o(1),
 \tag{FPK.29}
\]
where $\mathcal C_{\chi,k}>0$ is given by the literal constant
formula (FPK.22) with every $C_\pm$ replaced by $C_{\pm,k}$,
including their occurrences in the critical blocks (FPK.19).
The off-critical Cauchy blocks (FPK.16) remain unchanged.
This proves the corresponding fixed-packet result for the
programme's original reference, with its entire original mass.
